\chapter{Least Squares and Regression}

In scientific experiments and data analysis, empirical measurements rarely fit
theoretical models perfectly. Even when physical laws predict a linear relationship
between input and output (such as Hooke's law for an ideal spring), experimental
noise, measurement errors, and uncontrolled variables ensure that measured data
points scatter around the ideal curve.

In the language of linear algebra, we have an overdetermined system
$\mat{A}\vec{x}=\vec{b}$ whose measurement vector $\vec{b}$ lies outside the
column space $\col(\mat{A})$. Gaussian elimination reports ``no exact
solution,'' yet we do not throw the model away. Instead, we look for the
parameter vector $\hat{\vec{x}}$ that makes the \emph{error}
$\vec{e} = \vec{b} - \mat{A}\vec{x}$ -- the gap between measurements and model
predictions -- as small as possible. Finding $\hat{\vec{x}}$ requires no
iterative search; it comes straight from a square linear system, the
\textbf{normal equations}:
\[
  \mat{A}\T\mat{A}\,\hat{\vec{x}} \;=\; \mat{A}\T\vec{b}.
\]
Here, ``normal'' means perpendicular: the error vector stands at right angles to
every column of $\mat{A}$. Geometrically, $\mat{A}\hat{\vec{x}}$ is the orthogonal
projection of $\vec{b}$ onto $\col(\mat{A})$. Least squares applies the geometry of
orthogonal projections to data analysis, statistics, and machine learning.

\begin{focusbox}
When $\mat{A}\vec{x}=\vec{b}$ has no solution, do not solve it -- \emph{approximate} it. Find the
$\hat{\vec{x}}$ that makes $\mat{A}\hat{\vec{x}}$ as close to $\vec{b}$ as possible, measuring
closeness by the squared length $\norm{\vec{b}-\mat{A}\vec{x}}^2$. That $\hat{\vec{x}}$ solves
the normal equations
\[
  \mat{A}\T\mat{A}\,\hat{\vec{x}} \;=\; \mat{A}\T\vec{b},
\]
and geometrically $\mat{A}\hat{\vec{x}}$ is the \textbf{projection of $\vec{b}$ onto the column
space of $\mat{A}$}.
\end{focusbox}

\section{Why \texorpdfstring{$\mat{A}\vec{x}=\vec{b}$}{Ax=b} usually has no solution}

Let us be concrete before we are general. Suppose you believe two quantities are related by a
straight line $y = c + d\,t$, and you collect three data points:
\[
  (t_1,y_1)=(0,6),\qquad (t_2,y_2)=(1,0),\qquad (t_3,y_3)=(2,0).
\]
If a single line $y=c+dt$ passed through all three, the unknowns $c,d$ would have to satisfy
\[
  \begin{aligned}
    c + d\cdot 0 &= 6,\\
    c + d\cdot 1 &= 0,\\
    c + d\cdot 2 &= 0.
  \end{aligned}
\]
In matrix form $\mat{A}\vec{x}=\vec{b}$ with
\[
  \mat{A}=\begin{bmatrix} 1 & 0\\ 1 & 1\\ 1 & 2\end{bmatrix},
  \qquad
  \vec{x}=\colvec{c\\d},
  \qquad
  \vec{b}=\colvec{6\\0\\0}.
\]
Three equations, two unknowns. The first equation forces $c=6$. The second then forces $d=-6$.
But the third now reads $6 + 2(-6) = -6 \ne 0$. Contradiction. No line passes through all three
points. The system is \textbf{inconsistent} -- meaning there is no choice of parameters that can
satisfy every equation at once.

\begin{intuitionbox}
The columns of $\mat{A}$ span a two-dimensional subspace (a plane) in $\R^3$:
the span of $\colvec{1\\1\\1}$ and $\colvec{0\\1\\2}$. Every linear combination
$\mat{A}\vec{x}=c\colvec{1\\1\\1} + d\colvec{0\\1\\2}$ lies on this plane.
Because the measurement vector $\vec{b}=\colvec{6\\0\\0}$ does not lie in this
plane, no choice of coefficients $c,d$ can produce $\vec{b}$, and the system
$\mat{A}\vec{x}=\vec{b}$ has no exact solution.
\end{intuitionbox}

\begin{figure}[H]
\centering
\begin{tikzpicture}[scale=1.15,>=Stealth]
  % a tilted plane (the span of the two columns) drawn as a parallelogram
  \fill[sourceblue!12] (-1.5,-0.6) -- (2.6,-0.6) -- (3.6,1.2) -- (-0.5,1.2) -- cycle;
  \draw[sourceblue!55,thick] (-1.5,-0.6) -- (2.6,-0.6) -- (3.6,1.2) -- (-0.5,1.2) -- cycle;
  \node[sourceblue!70!black] at (3.1,1.0){$\col(\mat{A})$};
  \coordinate (O) at (0.0,0.0);
  \fill (O) circle (1.5pt) node[below left, font=\scriptsize]{$\vec{0}$};
  % the two column vectors living in the plane, pointing in clearly separated directions
  \draw[->,sourceblue!75!black,thick] (O)--(1.4,-0.15)
        node[below right=3pt,scale=0.85]{$\colvec{1\\1\\1}$};
  \draw[->,sourceblue!75!black,thick] (O)--(-0.35,0.75)
        node[above left=3pt,scale=0.85]{$\colvec{0\\1\\2}$};
  % b floating off the plane
  \coordinate (B) at (1.5,2.1);
  \draw[->,vecblue,very thick] (O)--(B) node[above right]{$\vec{b}=\colvec{6\\0\\0}$};
  \fill (B) circle (1.4pt);
  % dashed reminder that b is off the sheet: vertical drop to foot in the plane
  \coordinate (Bfoot) at (1.5,0.35);
  \draw[vecred,dashed,thick] (B)--(Bfoot);
  \fill[vecred] (Bfoot) circle (1.2pt);
  \node[spacegray,scale=0.85,right=4pt] at (1.5,1.25){off the plane};
\end{tikzpicture}
\caption{The two columns of $\mat{A}$, namely $(1,1,1)$ and $(0,1,2)$, span a flat plane through
the origin in $\R^3$. Every combination $\mat{A}\vec{x}=c(1,1,1)+d(0,1,2)$ lands somewhere on that
sheet. The data vector $\vec{b}=(6,0,0)$ points off the sheet, so no choice of $c,d$ reaches it:
the system $\mat{A}\vec{x}=\vec{b}$ is inconsistent.}
\label{ch19:plane-off}
\end{figure}

Figure~\ref{ch19:plane-off} draws exactly this: two columns spanning a sheet, and $\vec{b}$
hovering off it. This is the normal situation, not a pathology. Whenever you have \emph{more
equations than unknowns} -- more data points than parameters, which is what you always want in science -- the
matrix $\mat{A}$ is tall and skinny ($m>n$), its column space is a thin subspace of the big
output space $\R^m$, and a measured $\vec{b}$ almost never lands inside it. Demanding an exact
solution is demanding that noisy data fall exactly on a model. It will not.

\begin{trapbox}
``No solution'' from elimination does not mean ``no answer.'' It means there is no \emph{exact}
answer, which for over-determined data is the expected state. An exact fit to
many noisy points would usually mean your model has so many parameters that it is memorizing noise
(a hazard known as \emph{overfitting}). Inconsistent systems are resolved by orthogonal projection.
\end{trapbox}

\begin{recoverybox}
\textbf{Pause and think.} Suppose you have $m$ data points and fit a polynomial of degree $m-1$
(for example, two points and a line $y=c_0+c_1 t$, or three points and a parabola
$y=c_0+c_1 t+c_2 t^2$). Does $\mat{A}\vec{x}=\vec{b}$ have a solution?
Yes, always! When the number of parameters $n$ equals the number of distinct data points $m$, the
matrix $\mat{A}$ is square and invertible, so an exact polynomial passes through every point with zero
error. Inconsistency and the need for least squares only begin when $m > n$---when you collect more
measurements than your model has parameters, so the data overconstrains the equations.
\end{recoverybox}

\section{What ``closest'' should mean: the error vector}

We give up on $\mat{A}\vec{x}=\vec{b}$ and instead ask: of all the vectors $\mat{A}\vec{x}$ we
\emph{can} produce (all of which lie in the column space), which one is nearest to $\vec{b}$?
Define the \textbf{residual} or \textbf{error vector}
\[
  \vec{e} \;=\; \vec{b}-\mat{A}\vec{x}.
\]
The residual is simply what remains unexplained: each component $e_i = b_i - (\mat{A}\vec{x})_i$
measures how far the model's prediction misses the actual measurement. We want to make this whole
vector small. But ``small'' in what sense?

The choice that makes everything work is to minimize the \textbf{squared length}
\[
  \norm{\vec{e}}^2 \;=\; \norm{\vec{b}-\mat{A}\vec{x}}^2
  \;=\; \sum_{i=1}^{m} \bigl(b_i - (\mat{A}\vec{x})_i\bigr)^2 .
\]
This is the sum of the squared misses, and minimizing it is what gives the method its name:
\textbf{least squares}. For the line-fitting problem, $(\mat{A}\vec{x})_i = c+dt_i$ is the height
of the line above $t_i$, so $\norm{\vec{e}}^2 = \sum_i (y_i - c - d t_i)^2$ is the total squared
vertical distance from the data points to the line. Notice that we measure \emph{vertical}
distances, not perpendicular distances to the line: in scientific experiments, the input settings
$t_i$ (time, temperature, load) are treated as controlled and known, whereas the output
measurements $y_i$ absorb all the random experimental noise.

\begin{walkbox}
On a walk, ask yourself why we minimize the sum of \emph{squared} misses
$\sum_i e_i^2$ rather than plain absolute misses $\sum_i \abs{e_i}$. Compare
two things: what happens when you try to take a derivative of $\abs{e_i}$ at
zero, and which of the two sums equals a squared Euclidean length
$\norm{\vec{e}}^2$ where right triangles and the Pythagorean theorem apply?
\end{walkbox}

The moral to carry forward: squaring the errors turns data fitting into a
Euclidean \emph{geometry} problem, and geometry is something we can solve in
closed form. That is what the next section does.

\section{The geometry: projection onto the column space}

Here is the picture that organizes the entire chapter. All the achievable outputs $\mat{A}\vec{x}$
form the column space $\col(\mat{A})$ -- a plane (or higher-dimensional flat) sitting inside
$\R^m$. The data $\vec{b}$ is a point off that plane. We want the point \emph{in} the plane that
is closest to $\vec{b}$.

You already know what that point is: drop a perpendicular from $\vec{b}$ to the plane. The foot of
that perpendicular is the projection $\vec{p}=\proj_{\col(\mat{A})}\vec{b}$. The closest reachable
vector is $\vec{p}$, and the $\hat{\vec{x}}$ we seek is the one that produces it:
$\mat{A}\hat{\vec{x}}=\vec{p}$.

\begin{figure}[H]
\centering
\begin{tikzpicture}[scale=1.15,>=Stealth]
  % Plane C(A)
  \coordinate (O) at (0,0);
  \fill[sourceblue!12] (-2.5,-1.1) -- (2.7,-0.7) -- (3.5,1.1) -- (-1.7,0.7) -- cycle;
  \draw[sourceblue!70,thick] (-2.5,-1.1) -- (2.7,-0.7) -- (3.5,1.1) -- (-1.7,0.7) -- cycle;
  \node[draw=sourceblue!80, fill=sourceblue!15, rounded corners=2pt, font=\scriptsize\bfseries, text=sourceblue!95!black]
    at (2.7,0.65) {$\col(\mat{A}) \subset \R^m$};

  % Origin
  \fill[black] (O) circle (1.8pt) node[below left=2pt, font=\scriptsize] {$\vec{0}$};

  % Basis columns of A
  \coordinate (A1) at (1.8,-0.3);
  \coordinate (A2) at (-0.8,0.5);
  \draw[->,very thick,vecblue] (O) -- (A1) node[below right=1pt, font=\scriptsize] {$\vec{a}_1$};
  \draw[->,very thick,vecblue] (O) -- (A2) node[below left=1pt, font=\scriptsize] {$\vec{a}_2$};

  % Projection p in C(A)
  \coordinate (P) at (1.2,0.35);
  \draw[->,very thick,targetgreen!85!black] (O) -- (P);
  \node[targetgreen!85!black, font=\scriptsize, below right=1pt] at (P) {$\vec{p} = \mat{A}\hat{\vec{x}}$};
  \fill[targetgreen!85!black] (P) circle (1.8pt);

  % Observed data b off the plane
  \coordinate (B) at (1.2,2.6);
  \draw[->,very thick,vecred] (O) -- (B) node[above right=1pt, font=\scriptsize] {$\vec{b}$ (unreachable)};
  \fill[vecred] (B) circle (1.8pt);

  % Error vector e = b - p
  \draw[->,line width=1.6pt,quotientorange!95!black] (P) -- (B)
    node[midway,right=3pt,font=\small] {$\vec{e} = \vec{b} - \mat{A}\hat{\vec{x}}$};

  % Right-angle marker at P
  \draw[quotientorange!80!black,thick] (1.2,0.6) -- (0.95,0.6) -- (0.95,0.35);

  % Annotation below the plane
  \node[font=\scriptsize, sourceblue!90!black, align=center] at (0.4,-1.55) {
    $\vec{e} \perp \col(\mat{A}) \iff \mat{A}\T\vec{e} = \vec{0} \iff \mat{A}\T\mat{A}\hat{\vec{x}} = \mat{A}\T\vec{b}$
  };
\end{tikzpicture}
\caption{The fundamental geometry of least-squares regression. The reachable outputs $\mat{A}\vec{x}$
fill the column-space plane $\col(\mat{A})$. The data vector $\vec{b}$ floats off the plane.
The closest reachable point is the orthogonal projection $\vec{p}=\mat{A}\hat{\vec{x}}$.
The residual error vector $\vec{e} = \vec{b} - \mat{A}\hat{\vec{x}}$ meets the plane at a right angle,
orthogonal to every column of $\mat{A}$. That geometric perpendicularity $\mat{A}\T\vec{e}=\vec{0}$
yields the normal equations.}
\label{ch19:projection-picture}
\end{figure}

The key fact is the right angle. The shortest segment from a point to a plane is the
perpendicular one; any other segment from $\vec{b}$ into the plane is the hypotenuse of a right
triangle and hence longer. So at the optimum, the error $\vec{e}$ is \textbf{orthogonal to the
column space}. That single geometric statement is the whole derivation, and we now turn it into
algebra.

\begin{recoverybox}
\textbf{Pause and think.} What happens to this picture if $\vec{b}$ already lies \emph{inside} the
column space $\col(\mat{A})$?
Then $\vec{b}$ is already in the plane, so dropping a perpendicular moves it nowhere: the projection is
$\vec{p} = \vec{b}$, and the error vector is $\vec{e} = \vec{0}$. The normal equations become
$\mat{A}\T\mat{A}\hat{\vec{x}} = \mat{A}\T\vec{b}$, which is automatically satisfied by the exact
solution $\mat{A}\vec{x} = \vec{b}$. Least squares never distorts a problem that was already solvable;
it generalizes exact linear algebra to noisy real-world data.
\end{recoverybox}

\section{Deriving the normal equations}

The error $\vec{e}=\vec{b}-\mat{A}\hat{\vec{x}}$ must be orthogonal to every vector in
$\col(\mat{A})$. But every vector in the column space is a combination of the columns of
$\mat{A}$, so it is enough that $\vec{e}$ be orthogonal to \emph{each column} of $\mat{A}$.
Orthogonality of $\vec{e}$ to column $j$ means the dot product (column $j)\cdot\vec{e}=0$. Stacking
all $n$ of these dot products is exactly the matrix product $\mat{A}\T\vec{e}$. So the condition is
\[
  \mat{A}\T\vec{e} \;=\; \vec{0}
  \qquad\Longleftrightarrow\qquad
  \mat{A}\T\bigl(\vec{b}-\mat{A}\hat{\vec{x}}\bigr) \;=\; \vec{0}.
\]
Distribute and move one term to the other side:
\[
  \boxed{\;\mat{A}\T\mat{A}\,\hat{\vec{x}} \;=\; \mat{A}\T\vec{b}.\;}
\]
These are the \textbf{normal equations}. The word ``normal'' here means ``perpendicular'' -- it
records the right angle in Figure~\ref{ch19:projection-picture}. We started with an unsolvable tall
system $\mat{A}\vec{x}=\vec{b}$ ($m$ equations, $n$ unknowns) and multiplied both sides by
$\mat{A}\T$ to get a \emph{square} $n\times n$ system that is solvable.

\begin{formalbox}
\textbf{Least-squares solution.} Let $\mat{A}$ be $m\times n$ with linearly independent columns.
Then $\mat{A}\T\mat{A}$ is invertible, and the unique minimizer of $\norm{\vec{b}-\mat{A}\vec{x}}^2$
is
\[
  \hat{\vec{x}} \;=\; (\mat{A}\T\mat{A})\inv \mat{A}\T\vec{b}.
\]
The projection of $\vec{b}$ onto $\col(\mat{A})$ is $\vec{p}=\mat{A}\hat{\vec{x}}
=\mat{A}(\mat{A}\T\mat{A})\inv\mat{A}\T\vec{b}$, and the \textbf{projection matrix} onto the column
space is $\mat{P}=\mat{A}(\mat{A}\T\mat{A})\inv\mat{A}\T$.
\end{formalbox}

\begin{historybox}
On New Year's Day in 1801, the Italian astronomer Giuseppe Piazzi looked through his telescope in
Palermo and spotted a faint, moving point of light: the dwarf planet Ceres (the largest asteroid in
the belt between Mars and Jupiter). For 41 days, Piazzi tracked Ceres across just nine degrees of
the night sky before it disappeared into the blinding glare of the Sun. Circular orbit
approximations failed completely because Ceres follows an \emph{eccentric elliptical orbit} (an
elongated oval rather than a neat circle), and forty-one days of noisy telescope angles covered
barely one percent of its multi-year orbit. Most astronomers feared Ceres was lost for good.

Hearing of the crisis, a 24-year-old Carl Friedrich Gauss took up the challenge. Gauss had worked
out the method of least squares privately in 1795 as an eighteen-year-old student. Using least squares to
fit an overdetermined system of non-linear Keplerian equations (planetary equations derived from
Johannes Kepler's laws of orbital motion) to Piazzi's noisy tracking points, Gauss calculated
where Ceres would emerge on the other side of the Sun. On December 31, 1801, exactly one year after
Piazzi's discovery, the German astronomer Franz Xaver von Zach pointed his telescope precisely where
Gauss had directed and rediscovered Ceres within half a degree of his prediction.

In 1805, the French mathematician Adrien-Marie Legendre published the first public description of
least squares in his appendix on comet orbits (\emph{Nouvelles méthodes pour la détermination des
orbites des comètes}). When Gauss finally published his own celestial mechanics treatise in 1809
(\emph{Theoria Motus Corporum Coelestium}) and claimed priority by stating he had used the method
since 1795, a sharp priority feud erupted. Legendre protested bitterly, complaining to colleagues
that Gauss was claiming credit for a discovery Legendre had published first.

Earlier, Pierre-Simon Laplace had experimented with minimizing the sum of raw absolute errors
($\sum |e_i|$, the $L_1$ norm, which treats small and large errors with equal proportionality).
However, $L_1$ minimization leads to equations with sharp corners that were computationally
intractable before modern linear programming (an optimization technique developed in the mid-twentieth
century). The genius of minimizing squared errors ($L_2$, the ordinary Euclidean distance) is that
differentiating a smooth quadratic produces linear equations---the normal equations
$\mat{A}\T\mat{A}\hat{\vec{x}} = \mat{A}\T\vec{b}$---solvable by standard elimination.

The method unlocked celestial mechanics, geodesy (mapping the true curvature of the Earth), and
modern observational science. Yet least squares carries an honest physical weakness: squaring errors
heavily penalizes extreme outliers, making the fit hypersensitive to corrupted measurements. Finding
estimators that resist outliers while retaining computational speed remains an active subject in robust
statistics.
\end{historybox}

Let us also derive the normal equations a second way, by calculus, so you trust them no matter
which language you think in. Write the squared error as a function of $\vec{x}$:
\[
  f(\vec{x}) \;=\; \norm{\vec{b}-\mat{A}\vec{x}}^2
  \;=\; (\vec{b}-\mat{A}\vec{x})\T(\vec{b}-\mat{A}\vec{x})
  \;=\; \vec{b}\T\vec{b} - 2\,\vec{x}\T\mat{A}\T\vec{b} + \vec{x}\T\mat{A}\T\mat{A}\,\vec{x}.
\]
This is a quadratic bowl in $\vec{x}$. Setting its gradient to zero (using
$\nabla_{\vec{x}}(\vec{x}\T\mat{M}\vec{x})=2\mat{M}\vec{x}$ for symmetric $\mat{M}$, and
$\nabla_{\vec{x}}(\vec{x}\T\vec{c})=\vec{c}$):
\[
  \nabla f(\vec{x}) \;=\; -2\,\mat{A}\T\vec{b} + 2\,\mat{A}\T\mat{A}\,\vec{x} \;=\; \vec{0}
  \qquad\Longrightarrow\qquad
  \mat{A}\T\mat{A}\,\vec{x} = \mat{A}\T\vec{b}.
\]
Same equations. Because $\mat{A}\T\mat{A}$ is \textbf{positive semidefinite} (meaning
$\vec{x}\T(\mat{A}\T\mat{A})\vec{x} = \norm{\mat{A}\vec{x}}^2 \ge 0$ for every $\vec{x}$, so it
never turns negative), the bowl opens \emph{upward} in every direction like a soup bowl, so this
critical point is a genuine unique minimum, not a maximum or saddle point.

\begin{figure}[H]
\centering
\begin{tikzpicture}[scale=1.1,>=Stealth]
  % Coordinate axes for parameter space (c, d)
  \draw[->,spacegray,thick] (-0.5,0) -- (5.2,0) node[right,font=\small,text=black]{$c$ (intercept)};
  \draw[->,spacegray,thick] (0,-0.5) -- (0,3.8) node[above,font=\small,text=black]{$d$ (slope)};
  
  % The unique minimum point x-hat = (c-hat, d-hat)
  \coordinate (Xhat) at (2.4,1.7);
  
  % Concentric elliptical contours of the quadratic bowl
  \draw[sourceblue!45,thick,rotate around={25:(Xhat)}] (Xhat) ellipse (1.20 and 0.70);
  \draw[sourceblue!70,thick,rotate around={25:(Xhat)}] (Xhat) ellipse (2.00 and 1.20);
  \draw[sourceblue!95,thick,rotate around={25:(Xhat)}] (Xhat) ellipse (2.80 and 1.70);
  
  % Label on one contour line
  \node[sourceblue!80!black,font=\scriptsize,fill=white,inner sep=1.5pt] at ($(Xhat)+(-1.4,1.4)$) {$f(\vec{x})=\text{constant}$};
  
  % Gradient arrows pointing uphill (away from the minimum)
  \draw[->,vecred,thick] ($(Xhat)+(1.35,0.85)$) -- ($(Xhat)+(1.95,1.20)$) node[midway,above right,font=\scriptsize]{$\nabla f$};
  \draw[->,vecred,thick] ($(Xhat)+(-1.35,-0.85)$) -- ($(Xhat)+(-1.95,-1.20)$);
  
  % Dashed projections to axes for c-hat and d-hat
  \draw[spacegray,dashed] (Xhat) -- (2.4,0) node[below,font=\scriptsize,text=black]{$\hat{c}$};
  \draw[spacegray,dashed] (Xhat) -- (0,1.7) node[left,font=\scriptsize,text=black]{$\hat{d}$};

  % Center marker and label
  \fill[targetgreen!85!black] (Xhat) circle (2.2pt);
  \node[targetgreen!85!black,font=\small,above right=4pt,fill=white,fill opacity=0.85,text opacity=1,inner sep=1.5pt,rounded corners=2pt] at (Xhat) {$\hat{\vec{x}}$ (minimum)};
\end{tikzpicture}
\caption{Parameter space $\R^n$: the calculus view. The squared error $f(\vec{x})=\norm{\vec{b}-\mat{A}\vec{x}}^2$
is a smooth quadratic bowl whose level curves are concentric ellipses centered at the optimal
parameter vector $\hat{\vec{x}}$. The gradient $\nabla f(\vec{x})$ points uphill (red arrows). Setting
$\nabla f(\vec{x})=\vec{0}$ finds the unique bottom of the bowl in a single step, yielding the exact
same normal equations $\mat{A}\T\mat{A}\hat{\vec{x}}=\mat{A}\T\vec{b}$ that geometry found by
dropping a perpendicular in $\R^m$.}
\label{ch19:error-bowl}
\end{figure}

\begin{connectionbox}
Two completely different roads -- pure geometry (drop a perpendicular) and pure calculus (set the
gradient of a quadratic to zero) -- arrive at the exact same matrix equation. This is not a
coincidence; it is the linear-algebra shadow of a theme that runs through all of optimization.
Minimizing a smooth function means finding where its gradient vanishes, and for a \emph{quadratic}
function the gradient is linear, so ``gradient $=0$'' is a linear system. Least squares is the
simplest nontrivial optimization problem in the world, and it is the one case where the optimum has
a clean closed form. Every fancier method -- gradient descent, Newton's method, the whole edifice
of convex optimization -- can be understood as ``what do you do when the bowl is not exactly
quadratic.'' Here it is, so you get the answer in one step.
\end{connectionbox}

\subsection*{Why \texorpdfstring{$\mat{A}\T\mat{A}$}{A transpose A} is invertible}

We claimed that independent columns make $\mat{A}\T\mat{A}$ invertible. Here is the short proof,
worth seeing because it explains exactly when least squares fails. Suppose
$\mat{A}\T\mat{A}\,\vec{x}=\vec{0}$. Dot both sides with $\vec{x}$:
\[
  0 \;=\; \vec{x}\T\mat{A}\T\mat{A}\,\vec{x} \;=\; (\mat{A}\vec{x})\T(\mat{A}\vec{x})
  \;=\; \norm{\mat{A}\vec{x}}^2.
\]
A squared length is zero only when the vector is zero, so $\mat{A}\vec{x}=\vec{0}$. But if the
columns of $\mat{A}$ are independent, the only $\vec{x}$ with $\mat{A}\vec{x}=\vec{0}$ is
$\vec{x}=\vec{0}$. Thus the nullspace of $\mat{A}\T\mat{A}$ is trivial, and a square matrix with
trivial nullspace is invertible. The computation also shows $\vec{x}\T\mat{A}\T\mat{A}\,\vec{x}
=\norm{\mat{A}\vec{x}}^2\ge 0$, which is the positive-semidefiniteness that made the bowl open
upward.

\begin{trapbox}
If the columns of $\mat{A}$ are \emph{dependent} -- say you accidentally fed in the same feature
twice, or measured temperature in both Celsius and Fahrenheit -- then $\mat{A}\T\mat{A}$ is
singular and $(\mat{A}\T\mat{A})\inv$ does not exist. The projection $\vec{p}$ is still perfectly
well defined (the closest point on the plane is unique), but $\hat{\vec{x}}$ is not unique: many
different parameter vectors produce the exact same best fit. In statistics this is called
\textbf{collinearity} (or multicollinearity) -- meaning two or more feature columns point along
nearly the same direction, carrying redundant information that makes it impossible for the math to
disentangle their individual contributions. The cure is either to drop redundant features or to
add a shrinkage penalty (ridge regression, below), which restores invertibility.
\end{trapbox}

\section{Worked example: fitting a line}

Let us finish the three-point problem from the start of the chapter. We had
\[
  \mat{A}=\begin{bmatrix} 1 & 0\\ 1 & 1\\ 1 & 2\end{bmatrix},
  \qquad
  \vec{b}=\colvec{6\\0\\0},
  \qquad
  \hat{\vec{x}}=\colvec{c\\d}.
\]
Build the two pieces of the normal equations. First $\mat{A}\T\mat{A}$:
\[
  \mat{A}\T\mat{A}
  = \begin{bmatrix} 1&1&1\\ 0&1&2\end{bmatrix}
    \begin{bmatrix} 1&0\\ 1&1\\ 1&2\end{bmatrix}
  = \begin{bmatrix} 3 & 3\\ 3 & 5\end{bmatrix}.
\]
The top-left entry $3$ is just the number of data points; the off-diagonal $3=\sum t_i$ and the
bottom-right $5=\sum t_i^2$. Now $\mat{A}\T\vec{b}$:
\[
  \mat{A}\T\vec{b}
  = \begin{bmatrix} 1&1&1\\ 0&1&2\end{bmatrix}\colvec{6\\0\\0}
  = \colvec{6\\0}
  = \colvec{\sum y_i \\ \sum t_i y_i}.
\]
So the normal equations are the tidy $2\times 2$ system
\[
  \begin{bmatrix} 3 & 3\\ 3 & 5\end{bmatrix}\colvec{c\\d} = \colvec{6\\0}.
\]
Solve it. From the first row $3c+3d=6$, i.e.\ $c+d=2$. From the second $3c+5d=0$. Subtract three
times the first from the second: $(3c+5d)-(3c+3d)=2d=0-6=-6$, so $d=-3$, and then $c=5$. The
best-fit line is
\[
  \boxed{\,y = 5 - 3t.\,}
\]
Let us check the fit. The line predicts $\hat{y}=(5,2,-1)$ at $t=(0,1,2)$, while the data were
$\vec{b}=(6,0,0)$. So the projection and error are
\[
  \vec{p}=\mat{A}\hat{\vec{x}}=\colvec{5\\2\\-1},
  \qquad
  \vec{e}=\vec{b}-\vec{p}=\colvec{6\\0\\0}-\colvec{5\\2\\-1}=\colvec{1\\-2\\1}.
\]
Two sanity checks confirm the geometry. First, the residual should be orthogonal to both columns
of $\mat{A}$. Against the all-ones column: $1-2+1=0$. Against the time column $(0,1,2)$:
$0-2+2=0$. Perpendicular to the plane, exactly as the picture demanded. Second, the squared error
is $\norm{\vec{e}}^2 = 1+4+1=6$; no other line does better.
Figure~\ref{ch19:line-fit} shows the line, the data, and these residuals.

\begin{figure}[H]
\centering
\begin{tikzpicture}[scale=1.1,>=Stealth]
  \begin{axis}[
    axis lines=middle, xlabel={$t$}, ylabel={$y$},
    xmin=-0.4,xmax=2.6, ymin=-1.8,ymax=6.8,
    width=10cm,height=6.4cm,
    xtick={0,1,2}, ytick={-1,0,2,5,6},
    xticklabel style={anchor=north east},
    enlargelimits=false, clip=false,
  ]
    % best fit line y=5-3t
    \addplot[sourceblue,very thick,domain=-0.2:2.4]{5-3*x};
    \node[sourceblue!70!black] at (axis cs:0.55,4.2){$y=5-3t$};
    % data points
    \addplot[only marks,vecblue,mark=*,mark size=2.4pt] coordinates {(0,6) (1,0) (2,0)};
    % predicted points on the line
    \addplot[only marks,targetgreen!70!black,mark=square*,mark size=2.0pt] coordinates {(0,5) (1,2) (2,-1)};
    % residual segments (vertical)
    \draw[vecred,thick] (axis cs:0,6)--(axis cs:0,5);
    \node[left=3pt,font=\scriptsize,vecred] at (axis cs:0,5.5){$e_1=+1$};
    \draw[vecred,thick] (axis cs:1,0)--(axis cs:1,2);
    \node[right=3pt,font=\scriptsize,vecred] at (axis cs:1,1.0){$e_2=-2$};
    \draw[vecred,thick] (axis cs:2,0)--(axis cs:2,-1);
    \node[right=3pt,font=\scriptsize,vecred] at (axis cs:2,-0.5){$e_3=+1$};
  \end{axis}
\end{tikzpicture}
\caption{Fitting $y=c+dt$ to the three points $(0,6),(1,0),(2,0)$. The blue line is the
least-squares best fit $y=5-3t$. Filled circles are the data $\vec{b}$; green squares are the
predictions $\vec{p}=\mat{A}\hat{\vec{x}}$ on the line; red segments are the residuals
$e_i=b_i-p_i=(1,-2,1)$. Least squares makes the \emph{sum of the squared} red lengths as small as
possible.}
\label{ch19:line-fit}
\end{figure}

\begin{labbox}
Do the whole thing yourself with a fresh tiny data set: points $(t,y)=(-1,1),(0,1),(1,3)$. Build
$\mat{A}=\begin{bmatrix}1&-1\\1&0\\1&1\end{bmatrix}$, $\vec{b}=\colvec{1\\1\\3}$, form
$\mat{A}\T\mat{A}=\begin{bmatrix}3&0\\0&2\end{bmatrix}$ and $\mat{A}\T\vec{b}=\colvec{5\\2}$, and
read off $\hat{\vec{x}}=\colvec{5/3\\1}$. The best line is $y=\tfrac53 + t$. Verify the residual
$\vec{e}=\colvec{1\\1\\3}-\colvec{2/3\\5/3\\8/3}=\colvec{1/3\\-2/3\\1/3}$ is orthogonal to both
columns. (Notice how the centered times $t=-1,0,1$ made $\mat{A}\T\mat{A}$ \emph{diagonal} -- the
next box explains why this orthogonality is such a gift.)
\end{labbox}

\begin{aestheticbox}
Look at what just happened to the bookkeeping. The matrix $\mat{A}\T\mat{A}$ for line fitting is
\[
  \begin{bmatrix} m & \sum t_i \\ \sum t_i & \sum t_i^2 \end{bmatrix},
  \qquad
  \mat{A}\T\vec{b}=\colvec{\sum y_i \\ \sum t_i y_i}.
\]
Every entry is a \emph{sum} you compute by sweeping through the data once. The grand projection
geometry of $\R^m$ collapses into four little sums. This is why least squares scales: fitting a
line to a billion points costs the same matrix as fitting it to three -- you only ever build a
$2\times2$ system. The dimension of the data ($m$) never enters the size of the equations you
solve, only the number of \emph{parameters} ($n$) does. That decoupling is what makes regression
practical on enormous data.
\end{aestheticbox}

\section{Fitting a parabola, and any other model linear in its parameters}

Nothing about the method needed the model to be a line. Suppose the data look curved and you want
the best parabola $y=c_0+c_1 t+c_2 t^2$. The unknowns are $c_0,c_1,c_2$, and the model
remains strictly \emph{linear in those unknowns} even though it is quadratic in $t$. Each data point gives
one equation $c_0+c_1 t_i + c_2 t_i^2 = y_i$, so
\[
  \mat{A}=\begin{bmatrix}
    1 & t_1 & t_1^2\\
    1 & t_2 & t_2^2\\
    \vdots & \vdots & \vdots\\
    1 & t_m & t_m^2
  \end{bmatrix},
  \qquad
  \vec{x}=\colvec{c_0\\c_1\\c_2},
  \qquad
  \vec{b}=\colvec{y_1\\\vdots\\y_m}.
\]
This tall matrix (where each row contains increasing powers $1, t_i, t_i^2$ of an input sample,
a structure known as a \textbf{Vandermonde matrix}) translates polynomial evaluation into matrix
multiplication. It goes straight into the same normal equations
$\mat{A}\T\mat{A}\,\hat{\vec{x}}=\mat{A}\T\vec{b}$. The recipe is identical; only the number of
columns changed.

\begin{example}\label{ch19:parabola}
Fit $y=c_0+c_1t+c_2t^2$ to the four points $(t,y)=(-1,1),(0,1),(1,0),(2,4)$. Build
\[
  \mat{A}=\begin{bmatrix} 1&-1&1\\ 1&0&0\\ 1&1&1\\ 1&2&4\end{bmatrix},
  \qquad
  \vec{b}=\colvec{1\\1\\0\\4}.
\]
Then
\[
  \mat{A}\T\mat{A}
  =\begin{bmatrix} 4 & 2 & 6\\ 2 & 6 & 8\\ 6 & 8 & 18\end{bmatrix},
  \qquad
  \mat{A}\T\vec{b}=\colvec{6\\7\\17}.
\]
(Each entry of $\mat{A}\T\mat{A}$ is again a sum of powers of $t_i$: the top-left $4=\sum t_i^0=m$,
the next diagonal $6=\sum t_i^2$, the corner $18=\sum t_i^4$, and so on. The entries of
$\mat{A}\T\vec{b}$ are $\sum y_i=6$, $\sum t_i y_i=7$, and $\sum t_i^2 y_i=17$.) Solving the
$3\times3$ system by elimination gives
\[
  \hat{\vec{x}}=\colvec{c_0\\c_1\\c_2}=\colvec{1/10\\ -1/5 \\ 1}
  \quad\Longrightarrow\quad
  y = 0.1 - 0.2\,t + t^2 .
\]
The leading coefficient came out to $1$, recovering an underlying $t^2$ curvature, while the small
constant and linear terms absorbed the scatter. The predictions are $\vec{p}=\mat{A}\hat{\vec{x}}
=\bigl(\tfrac{13}{10},\tfrac{1}{10},\tfrac{9}{10},\tfrac{37}{10}\bigr)$, so the residual is
$\vec{e}=\vec{b}-\vec{p}=\bigl(-\tfrac{3}{10},\tfrac{9}{10},-\tfrac{9}{10},\tfrac{3}{10}\bigr)$,
which you can confirm is orthogonal to all three columns of $\mat{A}$ (each dot product is $0$).
Figure~\ref{ch19:parabola-fit} plots the data, the curve, and these residuals.
The fit is a true parabola, not a line, yet the normal equations stayed linear throughout.
\end{example}

\begin{figure}[H]
\centering
\begin{tikzpicture}[scale=1.1,>=Stealth]
  \begin{axis}[
    axis lines=middle, xlabel={$t$}, ylabel={$y$},
    xmin=-1.6,xmax=2.6, ymin=-0.6,ymax=4.6,
    width=10cm,height=6.4cm,
    xtick={-1,0,1,2}, ytick={0,1,2,3,4},
    enlargelimits=false, clip=false,
  ]
    % best-fit parabola y = 0.1 - 0.2 t + t^2
    \addplot[sourceblue,very thick,domain=-1.4:2.3,samples=80]{0.1-0.2*x+x^2};
    \node[sourceblue!70!black] at (axis cs:-0.95,3.7){$y=0.1-0.2t+t^2$};
    % data points
    \addplot[only marks,vecblue,mark=*,mark size=2.4pt] coordinates {(-1,1) (0,1) (1,0) (2,4)};
    % predicted points on the curve
    \addplot[only marks,targetgreen!70!black,mark=square*,mark size=2.0pt]
      coordinates {(-1,1.3) (0,0.1) (1,0.9) (2,3.7)};
    % residual segments (vertical)
    \draw[vecred,thick] (axis cs:-1,1)--(axis cs:-1,1.3);
    \draw[vecred,thick] (axis cs:0,1)--(axis cs:0,0.1);
    \draw[vecred,thick] (axis cs:1,0)--(axis cs:1,0.9);
    \draw[vecred,thick] (axis cs:2,4)--(axis cs:2,3.7);
  \end{axis}
\end{tikzpicture}
\caption{The least-squares parabola of Example~\ref{ch19:parabola}. Filled circles are the data
$\vec{b}=(1,1,0,4)$; green squares are the predictions $\vec{p}=\mat{A}\hat{\vec{x}}
=(\tfrac{13}{10},\tfrac{1}{10},\tfrac{9}{10},\tfrac{37}{10})$ on the curve; red segments are the
residuals. Exactly the same normal-equation machine that fit a line now fits a curve -- only the
number of feature columns changed.}
\label{ch19:parabola-fit}
\end{figure}

\begin{focusbox}
``Linear regression'' does not mean ``fit a straight line.'' It means \textbf{linear in the
parameters}. You may use any features of your inputs you like -- powers $t,t^2,t^3$, products of
two variables, $\sin t$, $\log t$, a hand-crafted indicator like ``is it a weekend'' -- and as long
as the unknown coefficients enter \emph{linearly}, the fit is a least-squares problem with the same
normal equations. The columns of $\mat{A}$ are your chosen features evaluated at each data point.
This single idea -- build a feature matrix, then project -- is most of practical machine learning
before neural networks.
\end{focusbox}

\section{Measuring the fit: residuals and \texorpdfstring{$R^2$}{R-squared}}

A best-fit line is only as good as the data allows. How good is it? The currency that answers this
is the residual vector $\vec{e}=\vec{b}-\mat{A}\hat{\vec{x}}$. Its squared length
$\norm{\vec{e}}^2=\sum e_i^2$ is the \textbf{residual sum of squares}, often written $\mathrm{SS_{res}}$.
Smaller is better, but it has units (squared $y$-units) and scale, so on its own a number like
``$6$'' means nothing. We want a unitless score that says ``how much of the variation in $y$ did
the model explain?''

That score is $R^2$, the \textbf{coefficient of determination}, and it falls straight out of a
Pythagorean theorem. Compare your model against the laziest possible model: predicting the constant
mean $\bar{y}=\tfrac1m\sum y_i$ for every point. The lazy model's error is the
\textbf{total sum of squares}
\[
  \mathrm{SS_{tot}} \;=\; \sum_{i=1}^m (y_i-\bar{y})^2,
\]
the raw variation in the data. Then define
\[
  \boxed{\;R^2 \;=\; 1 - \dfrac{\mathrm{SS_{res}}}{\mathrm{SS_{tot}}}
  \;=\; 1 - \dfrac{\norm{\vec{b}-\mat{A}\hat{\vec{x}}}^2}{\norm{\vec{b}-\bar{y}\,\mathbf{1}}^2}.\;}
\]
If the model fits perfectly, $\mathrm{SS_{res}}=0$ and $R^2=1$. If the model does no better than
predicting the mean, $\mathrm{SS_{res}}=\mathrm{SS_{tot}}$ and $R^2=0$. In between, $R^2$ is the
fraction of the variance the model captured.

\begin{intuitionbox}
$R^2$ compares two orthogonal distances. Total variation $\mathrm{SS_{tot}}$
measures how far the data spread from the horizontal line at the sample mean
$\bar{y}$; residual variation $\mathrm{SS_{res}}$ measures how far the data spread
from the fitted regression line. Because the model contains the constant mean as a
special case (when slope $d=0$), we always have $\mathrm{SS_{res}}\le\mathrm{SS_{tot}}$.
An $R^2$ of $0.9$ means that the linear model accounts for $90\%$ of the total
sample variance in the response variable, leaving $10\%$ in the residual error.
\end{intuitionbox}

\begin{recoverybox}
\textbf{Pause and think.} Can $R^2$ ever be negative?
Under the standard definition $R^2 = 1 - \mathrm{SS_{res}}/\mathrm{SS_{tot}}$, yes! If you fit a
model without an intercept (forcing the line through the origin) and the true data sit far above the
origin, the fitted line can make larger errors than the simple horizontal average $\bar{y}$. Then
$\mathrm{SS_{res}} > \mathrm{SS_{tot}}$, and $R^2$ drops below zero. A negative $R^2$ indicates
that the model fits the data worse than the simple sample mean $\bar{y}$.
\end{recoverybox}

\begin{example}\label{ch19:rsquared}
For our line fit $y=5-3t$ we had $\vec{b}=(6,0,0)$ and residual $\vec{e}=(1,-2,1)$, so
$\mathrm{SS_{res}}=\norm{\vec{e}}^2=6$. The mean is $\bar{y}=(6+0+0)/3=2$, so the deviations from
the mean are $(6-2,\,0-2,\,0-2)=(4,-2,-2)$ and
\[
  \mathrm{SS_{tot}} = 4^2+(-2)^2+(-2)^2 = 16+4+4 = 24.
\]
Therefore
\[
  R^2 = 1-\frac{6}{24} = 1-\frac14 = 0.75 .
\]
The line explains $75\%$ of the variation in this (tiny, noisy) data set, leaving a quarter
unexplained. With only three points and clear scatter, that is about what you would expect.
\end{example}

\begin{figure}[H]
\centering
\begin{tikzpicture}[scale=1.0,>=Stealth]
  % the column-space plane (the fitting subspace, containing the all-ones direction)
  \fill[sourceblue!12] (-1.7,-0.5) -- (2.5,-0.5) -- (3.3,1.0) -- (-0.9,1.0) -- cycle;
  \draw[sourceblue!55,thick] (-1.7,-0.5) -- (2.5,-0.5) -- (3.3,1.0) -- (-0.9,1.0) -- cycle;
  \node[sourceblue!70!black,font=\small] at (-0.9,0.7){\small $\col(\mat{A})$};
  \coordinate (O) at (0.0,0.1);
  % mean vector ybar*1 in the plane
  \coordinate (M) at (-0.5,0.0);
  % fit p in the plane
  \coordinate (P) at (1.7,0.35);
  % data b above the plane
  \coordinate (B) at (2.0,2.1);
  % the three legs
  \draw[->,spacegray,thick] (M)--(P) node[midway,below]{$\sqrt{\mathrm{SS_{model}}}$};
  \draw[->,vecred,very thick] (P)--(B) node[midway,right=2pt]{$\sqrt{\mathrm{SS_{res}}}\;=\;\norm{\vec{e}}$};
  \draw[->,vecblue,very thick] (M)--(B) node[midway,left=2pt,fill=white,inner sep=1pt]{$\sqrt{\mathrm{SS_{tot}}}$};
  % right angle at P aligned with the triangle legs
  \draw[vecred,thick] (1.50,0.32) -- (1.53,0.52) -- (1.73,0.55);
  \node[below left] at (M){$\bar{y}\,\mathbf{1}$};
  \node[below right] at (P){$\vec{p}=\mat{A}\hat{\vec{x}}$};
  \node[above right] at (B){$\vec{b}$};
  \fill (M) circle (1.2pt); \fill (P) circle (1.2pt); \fill (B) circle (1.2pt);
\end{tikzpicture}
\caption{$R^2$ is one right triangle. Measured from the mean vector $\bar{y}\,\mathbf{1}$, the data
$\vec{b}$ reaches up a distance $\sqrt{\mathrm{SS_{tot}}}$ (the hypotenuse). The fit $\vec{p}$ stays
in the column space, a leg of length $\sqrt{\mathrm{SS_{model}}}$; the residual $\vec{e}=\vec{b}
-\vec{p}$ is the perpendicular leg of length $\sqrt{\mathrm{SS_{res}}}$. Pythagoras gives
$\mathrm{SS_{tot}}=\mathrm{SS_{model}}+\mathrm{SS_{res}}$, and $R^2=1-\mathrm{SS_{res}}/
\mathrm{SS_{tot}}$ is the squared fraction of the hypotenuse used up by the model leg.}
\label{ch19:rsquared-triangle}
\end{figure}

\begin{connectionbox}
The Pythagorean theorem is hiding in plain sight. Because the residual $\vec{e}$ is orthogonal to
the column space and the fit $\vec{p}=\mat{A}\hat{\vec{x}}$ lives in it, $\vec{b}=\vec{p}+\vec{e}$
splits into perpendicular pieces, and $\norm{\vec{b}}^2=\norm{\vec{p}}^2+\norm{\vec{e}}^2$.
Subtracting the mean (which the all-ones column lets the fit reproduce) turns this into
$\mathrm{SS_{tot}}=\mathrm{SS_{model}}+\mathrm{SS_{res}}$, the famous \emph{analysis of variance}
(ANOVA) decomposition drawn in Figure~\ref{ch19:rsquared-triangle}. The correlation coefficient
$r$ is $\pm\sqrt{R^2}$, which equals the cosine of the angle between the centered data and the
centered predictor. This geometric decomposition underlies linear statistical modeling.
\end{connectionbox}

\begin{trapbox}
A high $R^2$ does \emph{not} mean your model is correct or useful. You can always raise $R^2$ by
adding more features (more columns to $\mat{A}$): in the extreme, fit $m$ points with an
$(m{-}1)$-degree polynomial and you get $R^2=1$ exactly -- the curve threads every point -- yet the
fit is a wild, useless wiggle that will predict the next data point terribly. This is
\textbf{overfitting} -- when a model memorizes random measurement noise rather than uncovering the
underlying relationship, destroying its ability to generalize. That is why practitioners report
\emph{adjusted} $R^2$, hold out separate test data, or penalize parameter size. $R^2$ measures fit to
the training sample you have, not truth, and never prediction on unseen data.
\end{trapbox}

\section{Least squares as the engine of statistics and machine learning}

Everything above was phrased as ``fit a curve,'' but the same three lines of algebra are the core
of supervised learning. Let us translate.

In machine learning you have $m$ training examples. Each example $i$ has a row of
\textbf{features} $\vec{a}_i\T = (a_{i1},\dots,a_{in})$ (the input attributes) and a numerical
\textbf{label} $b_i$ (the target value to predict). Stack the feature rows into the \textbf{design
matrix} $\mat{A}$ (the table of all inputs, with one row per example and one column per feature)
and the target labels into $\vec{b}$. A \textbf{linear model} predicts $\hat{b}_i=\vec{a}_i\T
\vec{x}$ using a weight vector $\vec{x}$ -- exactly $\mat{A}\vec{x}$. Training the model to minimize
squared prediction error is \emph{precisely} least squares, and the optimal weights are
\[
  \hat{\vec{x}} = (\mat{A}\T\mat{A})\inv\mat{A}\T\vec{b}.
\]
The fitted weights tell you how much each feature pushes the prediction; the residuals tell you
where the model is wrong; $R^2$ tells you how much signal you captured. This is \textbf{linear
regression}, the first model in every machine-learning course and still, today, the model that
runs quietly underneath an enormous fraction of deployed prediction systems because it is fast,
interpretable, and hard to break.

\begin{applicationbox}
\textbf{Why regression is everywhere.} A bank scoring loan applications stacks each applicant's
features (income, debt ratio, credit age) as a row of $\mat{A}$ and past repayment outcomes as
$\vec{b}$; the weights $\hat{\vec{x}}$ become the scorecard. A climate model fits temperature
against $\mathrm{CO}_2$, solar cycle, and aerosol features. A sports analyst predicts a player's
value from box-score features. In all of them the computational heart is the same projection onto
a column space.

Three engineering notes from real practice:
\begin{itemize}
  \item \textbf{Never invert $\mat{A}\T\mat{A}$ directly} on a real problem. Forming
  $\mat{A}\T\mat{A}$ squares the matrix's \textbf{condition number} $\kappa(\mat{A})$ (the ratio of
  its largest stretch factor to its smallest, measuring how severely round-off errors in the data
  get magnified in the answer). Instead, practical software solves least squares via the
  $\mat{Q}\mat{R}$ factorization (the subject of the next chapter) or the SVD, which work directly
  on $\mat{A}$ and are numerically stable. The normal equations are how you \emph{understand} least
  squares, not how you compute it in floating point.
  \item \textbf{Collinear features} (nearly dependent columns) make $\mat{A}\T\mat{A}$ nearly
  singular, so tiny data changes swing the weights wildly. The standard fix is \textbf{ridge
  regression}: minimize $\norm{\vec{b}-\mat{A}\vec{x}}^2 + \lambda\norm{\vec{x}}^2$, whose normal
  equations are $(\mat{A}\T\mat{A}+\lambda\mat{I})\hat{\vec{x}}=\mat{A}\T\vec{b}$. The added
  $\lambda\mat{I}$ guarantees invertibility and shrinks the weights toward zero.
  \item \textbf{Gradient descent} solves the very same problem one step at a time when $\mat{A}$ is
  too large to factor: it walks downhill on the bowl $f(\vec{x})=\norm{\vec{b}-\mat{A}\vec{x}}^2$
  (Figure~\ref{ch19:error-bowl}), and because the bowl is convex it cannot get stuck. The
  least-squares bowl is the training ground on which every optimizer in deep learning was first
  tested.
\end{itemize}
\end{applicationbox}

\begin{strangbox}
Strang frames least squares as the meeting point of the whole first half of linear
algebra: it combines the column space, orthogonality (the right angle of the error),
and projections (the matrix $\mat{P}$) into the single normal system
$\mat{A}\T\mat{A}\hat{\vec{x}}=\mat{A}\T\vec{b}$. Before Strang, keep the geometric
picture in mind: data vector off the subspace, drop a perpendicular. After Strang,
notice how this same equation reappears across advanced topics -- as the Moore--Penrose
pseudoinverse $\mat{A}^{+}=(\mat{A}\T\mat{A})\inv\mat{A}\T$, as the minimum-norm
solution via the SVD, and as the linear layer of a regression model. Gilbert Strang
calls $\mat{A}\T\mat{A}$ the fundamental matrix of applied mathematics, with least
squares providing its primary motivation.
\end{strangbox}

\section{The projection matrix, revisited}

Pulling the threads together, the fitted values are $\vec{p}=\mat{A}\hat{\vec{x}}=\mat{P}\vec{b}$
where
\[
  \mat{P} \;=\; \mat{A}(\mat{A}\T\mat{A})\inv\mat{A}\T.
\]
This $\mat{P}$ is the projection matrix onto $\col(\mat{A})$, and it carries two signature
properties that you can verify directly from its formula:
\[
  \mat{P}\T=\mat{P}\quad(\text{symmetric}),
  \qquad
  \mat{P}^2=\mat{P}\quad(\text{idempotent: projecting twice changes nothing}).
\]
The symmetry check: $\mat{P}\T=\bigl(\mat{A}(\mat{A}\T\mat{A})\inv\mat{A}\T\bigr)\T
=\mat{A}\bigl((\mat{A}\T\mat{A})\inv\bigr)\T\mat{A}\T=\mat{A}(\mat{A}\T\mat{A})\inv\mat{A}\T=\mat{P}$,
using that $\mat{A}\T\mat{A}$ is symmetric so its inverse is too. The idempotence check telescopes
beautifully:
\[
  \mat{P}^2 = \mat{A}\underbrace{(\mat{A}\T\mat{A})\inv\mat{A}\T\mat{A}}_{=\;\mat{I}}
  (\mat{A}\T\mat{A})\inv\mat{A}\T
  = \mat{A}(\mat{A}\T\mat{A})\inv\mat{A}\T = \mat{P}.
\]
The middle $\mat{A}\T\mat{A}$ cancels its own inverse, leaving $\mat{P}$ untouched. Geometrically:
once you have landed on the plane, projecting again does nothing, because you are already there.

\begin{example}\label{ch19:Pmatrix}
For our line-fitting $\mat{A}=\begin{bmatrix}1&0\\1&1\\1&2\end{bmatrix}$ we found
$\mat{A}\T\mat{A}=\begin{bmatrix}3&3\\3&5\end{bmatrix}$, whose inverse is
$\frac{1}{6}\begin{bmatrix}5&-3\\-3&3\end{bmatrix}$. Multiplying out
$\mat{P}=\mat{A}(\mat{A}\T\mat{A})\inv\mat{A}\T$ gives
\[
  \mat{P}=\frac16\begin{bmatrix} 5 & 2 & -1\\ 2 & 2 & 2\\ -1 & 2 & 5\end{bmatrix}.
\]
Check it against the data: $\mat{P}\vec{b}=\frac16\begin{bmatrix}5&2&-1\\2&2&2\\-1&2&5\end{bmatrix}
\colvec{6\\0\\0}=\frac16\colvec{30\\12\\-6}=\colvec{5\\2\\-1}=\vec{p}$, matching the projection we
computed by solving the normal equations. And $\mat{P}$'s rows sum to $1$ (it preserves the
constant vector, since the constant lives in the column space), while $\tr\mat{P}
=\tfrac{5+2+5}{6}=2$ equals the rank -- the dimension of the plane we projected onto.
\end{example}

\begin{remark}
The complementary matrix $\mat{I}-\mat{P}$ projects onto the orthogonal complement of the column
space $(\col(\mat{A}))^\perp = \nullsp(\mat{A}\T)$ (the subspace of all vectors perpendicular to
every column of $\mat{A}$, known as the left nullspace), and it sends $\vec{b}$ directly to the
residual: $(\mat{I}-\mat{P})\vec{b}=\vec{b}-\vec{p}=\vec{e}$. So $\mat{P}$ splits every data vector
into ``signal that the model can explain'' ($\mat{P}\vec{b}$, in the column space) plus ``noise it
cannot'' ($\vec{e}$, perpendicular to it). That decomposition $\vec{b}=\mat{P}\vec{b}
+(\mat{I}-\mat{P})\vec{b}$ is the whole of least squares in one line.
\end{remark}

\section{Summary}

\begin{focusbox}
\textbf{Summary of core concepts.} When $\mat{A}\vec{x}=\vec{b}$ has no solution because
$\vec{b}$ lies outside the column space, project: the best approximate solution $\hat{\vec{x}}$
makes the error $\vec{e}=\vec{b}-\mat{A}\hat{\vec{x}}$ perpendicular to $\col(\mat{A})$,
which gives the normal equations
\[
  \mat{A}\T\mat{A}\,\hat{\vec{x}}=\mat{A}\T\vec{b},
  \qquad
  \hat{\vec{x}}=(\mat{A}\T\mat{A})\inv\mat{A}\T\vec{b}
  \ \text{ when columns are independent.}
\]
The fit is $\vec{p}=\mat{P}\vec{b}$ with $\mat{P}=\mat{A}(\mat{A}\T\mat{A})\inv\mat{A}\T$
symmetric and idempotent. To fit a line, parabola, or any model linear in its parameters, put your
features in the columns of $\mat{A}$ and solve the normal equations. The coefficient of determination
is $R^2=1-\mathrm{SS_{res}}/\mathrm{SS_{tot}}$. Linear regression is the direct application of
orthogonal projections to observational data.
\end{focusbox}

\begin{connectionbox}
Every time you fit a model to noisy data, you are using least squares to drop a perpendicular from
your measurements onto the subspace of achievable predictions. By replacing impossible exact
equations with orthogonal projections, we extract the best available model in any dimension.

Yet solving the normal equations $\mat{A}\T\mat{A}\hat{\vec{x}}=\mat{A}\T\vec{b}$ directly carries
a hidden numerical hazard. Forming the product $\mat{A}\T\mat{A}$ squares the matrix's
\textbf{condition number} $\kappa(\mat{A})$---the ratio of its greatest stretch factor to its smallest,
which measures how severely round-off errors get amplified:
\[
  \kappa(\mat{A}\T\mat{A}) \;=\; \bigl(\kappa(\mat{A})\bigr)^2.
\]
If two features in $\mat{A}$ are nearly collinear, $\kappa(\mat{A})$ might be $10^4$. Squaring it
yields $10^8$, which on standard computer hardware wipes out eight decimal places of precision,
turning an accurate scientific model into pure numerical noise.

Can we compute the least-squares projection \emph{without} ever forming $\mat{A}\T\mat{A}$? In
Chapter~20, we will see that the answer is yes: by replacing crooked columns with an orthonormal
basis through the \textbf{Gram--Schmidt process}, we factor $\mat{A} = \mat{Q}\mat{R}$. With
orthonormal columns $\mat{Q}$, the matrix product $\mat{Q}\T\mat{Q} = \mat{I}$ becomes the identity
matrix, avoiding $\mat{A}\T\mat{A}$ altogether and solving least squares with clean numerical
stability.
\end{connectionbox}

\section*{Exercises}

\begin{exercise}\label{ch19:ex-mental-twopoints}
\textbf{(Mental check: exact fit vs approximation.)} Suppose you fit a straight line $y=c+dt$ to
just two distinct points $(t_1,y_1)=(1,3)$ and $(t_2,y_2)=(3,7)$. Before writing down any matrix:
what is the residual vector $\vec{e}$? What is $R^2$? Explain why fitting $n$ independent parameters
to $n$ data points never requires least squares.
\end{exercise}

\begin{exercise}\label{ch19:ex-mental-colspace}
\textbf{(Mental check: data already in the column space.)} Suppose a vector $\vec{b}$ already lies
in the column space of $\mat{A}$, so that $\mat{A}\vec{x}_0 = \vec{b}$ has an exact solution. What
is the projection $\vec{p}=\mat{P}\vec{b}$? What does the residual vector $\vec{e}$ equal? Verify
that the normal equations $\mat{A}\T\mat{A}\hat{\vec{x}}=\mat{A}\T\vec{b}$ are satisfied by
$\hat{\vec{x}}=\vec{x}_0$.
\end{exercise}

\begin{exercise}\label{ch19:ex-line}
Fit a line $y=c+dt$ by least squares to the four points $(0,1),(1,3),(2,4),(3,4)$. Form
$\mat{A}\T\mat{A}$ and $\mat{A}\T\vec{b}$, solve the $2\times2$ normal equations, and report the
best-fit line. Then compute the residual vector and verify it is orthogonal to both columns of
$\mat{A}$.
\end{exercise}

\begin{exercise}\label{ch19:ex-rsquared}
For your fit in Exercise~\ref{ch19:ex-line}, compute $\mathrm{SS_{res}}$, $\mathrm{SS_{tot}}$
(using the mean $\bar{y}$), and $R^2$. In one sentence, say what the resulting $R^2$ means about
the data.
\end{exercise}

\begin{exercise}\label{ch19:ex-horizontal}
Show that fitting the constant model $y=c$ (a horizontal line, so $\mat{A}$ is the single all-ones
column $\mathbf{1}$) by least squares gives $\hat{c}=\bar{y}$, the mean of the data. Conclude that
the mean is the least-squares ``best single number'' summarizing a data set, and explain why this
makes $R^2=0$ for the constant model.
\end{exercise}

\begin{exercise}\label{ch19:ex-parabola}
Fit a parabola $y=c_0+c_1 t+c_2 t^2$ to the three points $(-1,1),(0,0),(1,1)$. (With three points
and three parameters the fit should be \emph{exact}; predict before computing that the residual
will be zero, then confirm it by solving and checking $R^2=1$.)
\end{exercise}

\begin{exercise}\label{ch19:ex-centering}
Suppose you shift all the times so they sum to zero (replace $t_i$ by $t_i-\bar{t}$). Show that
then $\mat{A}\T\mat{A}=\begin{bmatrix} m & 0\\ 0 & \sum t_i^2\end{bmatrix}$ is diagonal, and hence
the best slope and intercept decouple into $\hat{d}=\dfrac{\sum t_i y_i}{\sum t_i^2}$ and
$\hat{c}=\bar{y}$. Explain geometrically why centering the data makes the two columns of $\mat{A}$
orthogonal.
\end{exercise}

\begin{exercise}\label{ch19:ex-P-properties}
Let $\mat{P}=\mat{A}(\mat{A}\T\mat{A})\inv\mat{A}\T$ for an $m\times n$ matrix $\mat{A}$ with
independent columns. Prove directly from the formula that (a) $\mat{P}\vec{p}=\vec{p}$ for any
$\vec{p}$ in the column space, and (b) $\mat{P}\vec{v}=\vec{0}$ for any $\vec{v}$ orthogonal to
every column of $\mat{A}$. Conclude that $\mat{P}$ acts as the identity on $\col(\mat{A})$ and as
zero on its orthogonal complement.
\end{exercise}

\begin{exercise}\label{ch19:ex-trace-rank}
Using the cyclic property $\tr(\mat{X}\mat{Y})=\tr(\mat{Y}\mat{X})$, show that $\tr\mat{P}=n$ for
the projection matrix of an $m\times n$ matrix with independent columns. Interpret this number
geometrically (what is the dimension of the space $\mat{P}$ projects onto?).
\end{exercise}

\begin{exercise}\label{ch19:ex-ridge}
\textbf{(Ridge regression.)} For scalar data, suppose $\mat{A}=\mathbf{1}$ (the all-ones column,
$m$ entries) and we minimize $\norm{\vec{b}-c\mathbf{1}}^2+\lambda c^2$. Derive the modified normal
equation $(m+\lambda)\hat{c}=\sum b_i$, so $\hat{c}=\dfrac{m\,\bar{y}}{m+\lambda}$. Show that as
$\lambda\to 0$ you recover the mean, and as $\lambda\to\infty$ the estimate shrinks to $0$.
Explain in words why a penalty $\lambda\norm{\vec{x}}^2$ ``shrinks'' the fit.
\end{exercise}

\begin{exercise}\label{ch19:ex-two-features}
\textbf{(Multiple regression.)} You want to predict $b$ from two features $u$ and $v$ with the
model $b\approx x_1 u + x_2 v$ (no intercept). The data are
$(u,v,b)=(1,0,1),(0,1,1),(1,1,3)$. Set up $\mat{A}=\begin{bmatrix}1&0\\0&1\\1&1\end{bmatrix}$,
$\vec{b}=\colvec{1\\1\\3}$, solve the normal equations for $\hat{\vec{x}}$, and interpret each
weight as the effect of one feature holding the other fixed.
\end{exercise}

\begin{exercise}\label{ch19:ex-why-squares}
\textbf{(A proof about minimizing.)} Let $\vec{p}$ be the projection of $\vec{b}$ onto a subspace
$S=\col(\mat{A})$, so $\vec{b}-\vec{p}\perp S$. Take \emph{any} other vector $\vec{q}\in S$ and
prove $\norm{\vec{b}-\vec{q}}^2=\norm{\vec{b}-\vec{p}}^2+\norm{\vec{p}-\vec{q}}^2$ (Pythagoras),
hence $\norm{\vec{b}-\vec{q}}\ge\norm{\vec{b}-\vec{p}}$ with equality only when $\vec{q}=\vec{p}$.
This is the rigorous statement that the projection is the unique closest point -- the foundation the
whole chapter rests on.
\end{exercise}
